\section{A positive four-supercharge enlargement}
\label{ext-section}

An enlargement of the field space provides a concrete way to obtain four
real supercharges without retaining the spectrum of a two-supercharge
model.  The construction below adds an angular coordinate to the real
Morse variable.  It retains the gamma function as a holomorphic period,
but its physical vacuum is a differential form on the enlarged target.
Adding two complex oscillators then produces a physical Euler-factor
metric in an explicitly specified parameter frame.  These statements
concern positive quantum theories and identifiable observables; they do
not identify their vacuum pairing with the complete Weil form.

\subsection{The cylinder and its supersymmetry algebra}
\label{ext-cylinder}

Let
\begin{equation}
 \mathcal X=\R\times(\R/2\pi\mathbb Z),\qquad
 Y=y+i\theta,\qquad
 \mathcal H=L^2\Omega^*(\mathcal X;\C),
 \label{ext-hilbert}
\end{equation}
with flat metric $dy^2+d\theta^2$ and the ordinary positive inner product
on complex differential forms.  Write $E_y,E_\theta$ for exterior
multiplication by $dy,d\theta$, and $I_y,I_\theta$ for their Hilbert
adjoints.  Thus $\{E_i,I_j\}=\delta_{ij}$, while the anticommutators of
two exterior multiplications or two contractions vanish.  Initially
take $a>0$ and define
\begin{equation}
 \mathcal W_a(Y)=\frac12(e^Y-aY),\qquad
 h=\re\mathcal W_a
   =\frac12(e^y\cos\theta-ay).
 \label{ext-superpotential}
\end{equation}
Although $\mathcal W_a$ has an additive period, $h$ is single valued for
real $a$.  In particular $\Delta h=h_{yy}+h_{\theta\theta}=0$.

On smooth compactly supported periodic forms define
\begin{align}
 Q_1&=E_y(\partial_y+h_y)+E_\theta(\partial_\theta+h_\theta),\notag\\
 Q_1^\dagger&=I_y(-\partial_y+h_y)
                    +I_\theta(-\partial_\theta+h_\theta),\notag\\
 Q_2&=-E_y(\partial_\theta-h_\theta)+E_\theta(\partial_y-h_y),\notag\\
 Q_2^\dagger&=I_y(\partial_\theta+h_\theta)
                    +I_\theta(-\partial_y-h_y).
 \label{ext-complex-charges}
\end{align}
Equivalently, $Q_1=d+dh\wedge$ and $Q_2=d^c-d^ch\wedge$, where
$d^c=E_\theta\partial_y-E_y\partial_\theta$.  All adjoints in
\eqref{ext-complex-charges} refer to \eqref{ext-hilbert}; an indefinite
fermionic pairing is not involved.

For a general smooth real $h$, the exterior relations give
\begin{align}
 Q_1^2=Q_2^2&=0,\qquad \{Q_1,Q_2^\dagger\}=0,\notag\\
 \{Q_1,Q_2\}&=-2(\Delta h)E_yE_\theta,\notag\\
 \{Q_1,Q_1^\dagger\}-\{Q_2,Q_2^\dagger\}
   &=2(\Delta h)(N-1),\qquad
 N=E_yI_y+E_\theta I_\theta.
 \label{ext-algebra-defects}
\end{align}
For example, the coefficient of $E_yE_\theta$ in the mixed nilpotent
anticommutator is
$[\partial_y+h_y,\partial_y-h_y]
+[\partial_\theta+h_\theta,\partial_\theta-h_\theta]
=-2\Delta h$.  The remaining identities follow by the same exterior
expansion.  Harmonicity therefore gives a common Hamiltonian
$H=\{Q_1,Q_1^\dagger\}=\{Q_2,Q_2^\dagger\}$ and four real charges
\begin{equation}
 \begin{split}
 R_1=Q_1+Q_1^\dagger,&\qquad R_2=i(Q_1-Q_1^\dagger),\\
 R_3=Q_2+Q_2^\dagger,&\qquad R_4=i(Q_2-Q_2^\dagger),
 \end{split}
 \qquad \{R_i,R_j\}=2\delta_{ij}H.
 \label{ext-real-algebra}
\end{equation}
Their common square is
\begin{equation}
 \begin{split}
 H={}&-\partial_y^2-\partial_\theta^2+\frac14|e^Y-a|^2\\
 &+e^y\bigl[\cos\theta(E_yI_y-E_\theta I_\theta)
             -\sin\theta(E_yI_\theta+E_\theta I_y)\bigr].
 \end{split}
 \label{ext-hamiltonian}
\end{equation}
This is an exponential bosonic interaction with a Hessian coupling to
the fermions.  In the corresponding normal-ordered coherent-state
description its Euclidean action is
\begin{equation}
 S_E=\int dt\left\{
 \frac14(\dot y^2+\dot\theta^2)+|\nabla h|^2
 +\bar c_i\dot c_i+2\bar c_i h_{ij}c_j\right\}.
 \label{ext-action}
\end{equation}
The normalization follows from the kinetic operator $-\Delta$ in
\eqref{ext-hamiltonian}; the scalar normal-ordering contribution
$-\Delta h$ vanishes.  Positivity is the operator statement $H\geq0$,
not positivity of every fermionic path-integral integrand.

\subsection{Closed operators and the physical vacuum}
\label{ext-domains-vacuum}

The displayed algebra has a closed realization.  Each $R_i$ is a
Dirac-type first-order operator with smooth Hermitian multiplication
potential on a complete flat manifold.  Choose compact-support
cutoffs $\chi_L$ tending to one with $|d\chi_L|\leq C/L$.  The
potential commutes with the cutoff, so
$\|[R_i,\chi_L]\|\leq C/L$.  A deficiency vector satisfying
$R_i^*\psi=\pm i\psi$ is locally regular by ellipticity.  Integration
by parts with $\chi_L^2\psi$, followed by taking the imaginary part,
gives
\begin{equation}
 \|\chi_L\psi\|^2
 \leq \frac{C}{L}\|\chi_L\psi\|\,\|\psi\|.
 \label{ext-cutoff}
\end{equation}
Consequently both deficiency spaces vanish.  This proves essential
self-adjointness on the stated common core.  The same argument applies
to every nonzero real linear combination of the $R_i$.

On that core the graph norms of the four charges coincide, by
\eqref{ext-real-algebra}.  Their closures therefore have a common
quadratic-form domain, and their squares define the same nonnegative
self-adjoint $H$.  Polarizing the identities for real linear
combinations proves the mixed algebra as a form identity.  Since each
closed $R_i$ commutes with the spectral resolution of $R_i^2=H$, all
products in \eqref{ext-real-algebra} are defined on $\operatorname{Dom}H$,
where the operator identities follow.  This argument also applies to
the product targets used below.  It is compatible with the
noncompact Witten-operator framework of \cite{DaiYan}; the cutoff
argument itself does not require the vacuum-counting theorem.

Restricting the function $h$ to $\theta=0$ gives
$(e^y-ay)/2$, but restricting a function is not a quantum reduction.
For example, the scalar real-line operator obtained from
$A=\partial_y+(e^y-a)/2$ is
\begin{equation}
 A^\dagger A=-\partial_y^2+\frac14(e^y-a)^2-\frac12e^y.
 \label{ext-old-morse}
\end{equation}
The scalar cylinder Hamiltonian has no negative Hessian term because
$\Delta h=0$.  In degree one its Hessian term has eigenvalues
$\pm e^y$, rather than $-e^y/2$.  A delta function at a fixed angle
is not a state of \eqref{ext-hilbert}.  Moreover, $e^{-h}$ grows
exponentially near $\theta=\pi$ as $y\to+\infty$, and $e^h$ grows
near $\theta=0$.  Neither supplies a normalizable scalar or top-form
vacuum.

\begin{proposition}
\label{ext-vacuum-count}
For every $a\in\C\setminus\{0\}$, the cylinder Hamiltonian with
periodic differential forms has exactly one complex normalizable
zero-energy state.  It is a one-form.  The assertion uses the
well-tame Morse inequalities of \cite{DaiYan} at large Witten
deformation, together with the continuation argument below; it does
not compute a vacuum norm in an arithmetic source frame.
\end{proposition}

\begin{proof}
First let $a>0$ and replace $h$ by $Th$, $T>0$.  There is exactly one
critical point, $(y,\theta)=(\log a,0)$, of Morse index one.  The
complete flat cylinder has bounded geometry, and
\begin{equation}
 \liminf_{|y|\to\infty}|dh|=\frac a2>0,
 \qquad
 \frac{|\operatorname{Hess}h|}{|dh|^2}\longrightarrow0.
 \label{ext-tame}
\end{equation}
At the negative end the Hessian tends to zero and $|dh|\to a/2$;
at the positive end $|dh|$ grows like $e^y/2$ while the Hessian grows
like $e^y$.  These are the well-tame hypotheses in \cite{DaiYan}.
Their strong Morse inequalities, including equality of the full
alternating sums, give graded Dirac index $-1$ for sufficiently
large $T$.  This step uses the analytic theorem, rather than an
assumed identification of critical points with exact quantum states.
It does not require the additional Thom--Smale comparison theorem.

The graded Dirac operator remains Fredholm for all $T>0$.  At
$y\to-\infty$ its square tends to
$-\partial_y^2-\partial_\theta^2+T^2a^2/4$; at the other end its
potential tends to $+\infty$, since $T^2|dh|^2$ dominates the Hessian.
These estimates are uniform on compact positive $T$ intervals and
give a positive lower bound outside a compact set.  Local elliptic
compactness then gives Fredholmness.  The estimate
$|\operatorname{Hess}h|\leq\varepsilon|dh|^2+C_\varepsilon$
also makes the graph domains equivalent on those intervals, so the
family is continuous in graph norm and its index is constant.
This is the required continuation from the large-deformation
Morse calculation \cite{WittenMorse,DaiYan}.

In degrees zero and two the squared operator is simply
$-\Delta+T^2|dh|^2$.  A zero state would have both vanishing weak
gradient and vanishing product with $|dh|$, and hence must be zero.
Thus the index $-1$ implies precisely one degree-one vacuum.

For complex $a$, the globally defined one-form $dh$ is described
explicitly below.  Its variation with $a$ is a bounded constant
one-form.  The negative-end gap becomes $T^2|a|^2/4$ and the
positive-end estimate persists.  Any path in $\C\setminus\{0\}$
therefore gives a Fredholm continuation from positive real $a$.
The even-degree kernels remain zero because $\Delta h=0$ locally.
The same count follows at $T=1$.  At $a=0$ the end gap closes, so
this argument makes no assertion there.
\end{proof}

Writing the vacuum as $\omega=P\,dy+Q\,d\theta$, its equations are
\begin{equation}
 \begin{split}
 \partial_yQ-\partial_\theta P+h_yQ-h_\theta P&=0,\\
 -\partial_yP-\partial_\theta Q+h_yP+h_\theta Q&=0.
 \end{split}
 \label{ext-vacuum-equations}
\end{equation}
These determine the physical state on the full cylinder.  No scalar
gamma wavefunction is substituted for their solution.

\subsection{Periodic parameters and the gamma period}
\label{ext-periodic-parameter}

For $a=a_R+ia_I$ one has locally
\begin{equation}
 h=\frac12(e^y\cos\theta-a_Ry+a_I\theta),\qquad
 h(y,\theta+2\pi)-h(y,\theta)=\pi a_I.
 \label{ext-period}
\end{equation}
Although $h$ need not be globally a function, its differential is a
smooth closed real one-form on the cylinder.  The supercharges use
this one-form and its complex-structure rotation, so remain globally
defined on periodic forms.  Their definition is not globally a
conjugation by a single-valued $e^h$.  Similarly,
$\mathcal W_a(Y+2\pi i)-\mathcal W_a(Y)=-\pi ia$, and the proposed
chiral insertion $\partial_a\mathcal W_a=-Y/2$ is multivalued.
The local holomorphic dependence of a superpotential thus does not
by itself supply an ordinary globally defined rank-one chiral
$tt^*$ problem.  Periodic Landau--Ginzburg theories with deck and
flavor data provide a relevant framework \cite{CGV}; any such extra
data must be included explicitly in a proposed arithmetic pairing.

The ordinary Berry connection of the specific one-vacuum realization
above is locally flat.  Indeed, $H$ has real coefficients for every
complex $a$, and its isolated one-dimensional kernel admits a real
normalized generator $\psi$ locally.  Reality and differentiation of
$\langle\psi,\psi\rangle=1$ imply
$\langle\psi,d\psi\rangle=0$.  A possible global sign holonomy does
not change this local conclusion.  Parameter connections with
additional insertions or flavor structure are separate constructions;
the distinction matters when using supersymmetric Berry equations
\cite{SonnerTong}.

There is nevertheless an exact holomorphic gamma period:
\begin{equation}
 \mathcal P_\gamma(a)
 =\int_{Y\in\R}e^{-2\mathcal W_a(Y)}\,dY
 =\int_0^\infty t^{a-1}e^{-t}\,dt
 =\Gamma(a),\qquad \re a>0.
 \label{ext-gamma-period}
\end{equation}
The factor two in the exponent is part of the definition.  For
complex $a$ the real path is a convergent relative cycle, not
necessarily a steepest-descent thimble.  Such periods belong to the
holomorphic boundary quantities of Landau--Ginzburg theory
\cite{CGV}.  Equation \eqref{ext-gamma-period} is not an evaluation
of the Hermitian norm of the physical one-form in
\eqref{ext-vacuum-equations}.

\subsection{Two complex oscillators and a physical Euler metric}
\label{ext-euler}

Introduce independent parameters $a,q$ with $\re a>0$, $|q|<1$,
put $M=1-q$, and enlarge the target to $\mathcal X\times\C^2$ with
its product flat metric.  The superpotential is
\begin{equation}
 \mathcal W_{a,q}(Y,U,V)
 =\frac12\left[e^Y-aY+\frac M2(U^2+V^2)\right].
 \label{ext-product-superpotential}
\end{equation}
This introduces two further complex bosons and their form-valued
fermions.  On a flat Kähler target the preceding construction extends
with $d^c$ defined by the complex structure.  The required Hessian
condition is pluriharmonicity, rather than merely a vanishing trace;
it holds for $h=\re\mathcal W_{a,q}$.  Alternatively, the graded
tensor product of the three two-dimensional complexes verifies the
four-charge algebra directly.  The common-core closure argument
above applies on this complete product target.  In particular,
\begin{equation}
 |\nabla h|^2=\frac14|e^Y-a|^2
               +\frac{|M|^2}{4}(|U|^2+|V|^2).
 \label{ext-product-potential}
\end{equation}

Since $\re M>0$, elementary integration on the real cycle gives
\begin{equation}
 \frac1{2\pi}\int_{\R^3}e^{-2\mathcal W_{a,q}}\,dY\,dU\,dV
 =\frac{\Gamma(a)}{1-q}.
 \label{ext-product-period}
\end{equation}
For a fixed prime $p$, setting $a=z/2$, $q=p^{-z}$ and specifying the
source multiplier $\pi^{-z/2}$ produces
$\pi^{-z/2}\Gamma(z/2)(1-p^{-z})^{-1}$ for $\re z>0$.
Neither this restriction of parameters nor its Fourier variable
$\im z$ identifies an arithmetic coordinate with physical Euclidean
time.

The oscillator vacuum metric is separately computable.  For one
field with $\mathcal W_U=MU^2/4$, write $M=r e^{i\phi}$ and
$U'=e^{i\phi/2}U=x'+iy'$.  Then
\begin{equation}
 h_U=\frac r4(x'^2-y'^2),\qquad
 \psi_M=e^{-r|U|^2/4}\,dy',\qquad
 \|\psi_M\|^2=\frac{2\pi}{r}.
 \label{ext-oscillator-vacuum}
\end{equation}
The form $\psi_M$ satisfies $Q_1\psi_M=Q_1^\dagger\psi_M=0$.
Equivalently, the scalar oscillator ground energy is $r$, cancelled
by the occupied negative Hessian direction, whose contribution is
$-r$.  Every other oscillator state has positive energy, proving
uniqueness.  Consequently the product theory also has a unique
physical vacuum for $aM\ne0$.

To make the parameter frame explicit, define
\begin{equation}
 \begin{split}
 s_M&=(2\pi)^{-1/2}e^{-i\phi/2}\psi_M\\
 &=\frac{e^{-r|U|^2/4}}{2i\sqrt{2\pi}}
       \left(dU-\frac{\bar M}{|M|}\,d\bar U\right).
 \end{split}
 \label{ext-holomorphic-frame}
\end{equation}
The second expression is single valued for $M\ne0$.  If $P_M$
denotes the physical ground-state projector, then
$P_M\partial_{\bar M}s_M=0$.  To verify this without identifying
different Hodge realizations, normalize the real state as
$\chi_M=\sqrt{r/(2\pi)}\psi_M$.  It has zero local Berry connection,
and $s_M=M^{-1/2}\chi_M$ on a local branch.  The coefficient is
holomorphic, which proves the assertion directly in the Hilbert
space of real differential forms complexified.  Thus
\begin{equation}
 \langle s_M,s_M\rangle=|M|^{-1},\qquad
 g_{UV}=|M|^{-2}=|1-q|^{-2}.
 \label{ext-physical-euler-metric}
\end{equation}
No identification with a twisted-Dolbeault identity class is needed
for this explicit calculation.

The coefficient in \eqref{ext-holomorphic-frame} fixes the source
normalization.  In the associated tensor frame the full physical
metric is $g_\gamma(a,\bar a)|1-q|^{-2}$, where the cylinder factor
has not been evaluated.  It is distinct from the period
\eqref{ext-product-period}.  Furthermore,
$\partial_M\partial_{\bar M}\log g_{UV}=0$ away from zero, and
multiplication of the two-field frame by $M$ gives unit metric.
Thus the denominator has meaning relative to the specified source;
positivity alone does not select it.  At $M=0$ oscillator confinement
and its normalizable vacuum disappear.  These facts provide a
positive physical benchmark, while leaving the arithmetic contact
and full pole pairing undetermined.

\subsection{An affine gamma--prime coupling fails two tests}
\label{ext-affine}

A direct attempt to couple the sectors replaces their constant mass
by $M+\lambda e^Y$:
\begin{equation}
 \mathcal W^\lambda_{a,q}
 =\frac12\left[e^Y-aY+
             \frac{M+\lambda e^Y}{2}(U^2+V^2)\right].
 \label{ext-affine-superpotential}
\end{equation}
For real $a,M,\lambda>0$ its real-cycle integral converges, and
integrating $U,V$ gives
\begin{equation}
 \begin{split}
 \mathcal P_\lambda(a,M)
 &=\int_0^\infty\frac{t^{a-1}e^{-t}}{M+\lambda t}\,dt
 <\frac{\Gamma(a)}{M},\\
 \left.\frac{d}{d\lambda}\mathcal P_\lambda\right|_{0+}
 &=-\frac{\Gamma(a+1)}{M^2}.
 \end{split}
 \label{ext-affine-period}
\end{equation}
The inequality and the right derivative are exact analytic checks
of the amplitude.  They are not numerical positivity bounds.

There is a separate obstruction in the full complex target.  At
$t_0=-M/\lambda\ne0$, the critical equations admit
\begin{equation}
 e^Y=t_0,\qquad
 U^2+V^2=\frac{2(a-t_0)}{\lambda t_0}.
 \label{ext-critical-quadric}
\end{equation}
This is a noncompact complex critical quadric: writing
$(U+iV)(U-iV)$ as the displayed constant exhibits points escaping
to infinity.  Thus a convergent real contour conceals nonisolated
classical vacua elsewhere in the physical field space.  The
isolated-vacuum hypotheses of a standard massive ground-state
construction cannot simply be assumed.  This observation does not
by itself prove a vanishing quantum spectral gap.

The period obstruction has a useful scoped extension.  Fix $q$ and
suppose a mass $m(t,q)$ is nonzero on $t>0$ and satisfies the
integrability assumptions for uniqueness of its Mellin transform.
If
\begin{equation}
 \int_0^\infty t^{a-1}e^{-t}m(t,q)^{-1}\,dt
       =\frac{\Gamma(a)}{1-q}
 \label{ext-mellin-test}
\end{equation}
holds throughout an open Mellin strip for independently variable
$a$, then $m(t,q)=1-q$ almost everywhere on that contour.  Indeed,
subtract the two integrands, fix a real part $\re a=b$ in the strip,
and put $t=e^x$.  The resulting $L^1$ function of $x$ has identically
zero Fourier transform, so vanishes almost everywhere.  This
excludes a nonconstant quadratic mass preserving the entire
factorized period within this class.  It does not exclude a match
only on the linked locus $q=p^{-2a}$, a different insertion or
boundary problem, or a nonquadratic prime theory.  Those alternatives
require additional structure beyond the positive product
Hamiltonian constructed here.
